% ****** Unified Dynamics — AIP Format ******
%
%  Compile: pdflatex -> bibtex -> pdflatex -> pdflatex
%
\documentclass[%
 aip,
 amsmath,amssymb,
 reprint,%
]{revtex4-1}

\usepackage{graphicx}% Include figure files
\usepackage{dcolumn}% Align table columns on decimal point
\usepackage{bm}% bold math

\usepackage[utf8]{inputenc}
\usepackage[T1]{fontenc}
\usepackage{mathptmx}
\usepackage{etoolbox}
\usepackage{amsthm}

% AIP email placement patch
\makeatletter
\def\@email#1#2{%
 \endgroup
 \patchcmd{\titleblock@produce}
  {\frontmatter@RRAPformat}
  {\frontmatter@RRAPformat{\produce@RRAP{*#1\href{mailto:#2}{#2}}}\frontmatter@RRAPformat}
  {}{}
}%
\makeatother

% Theorem environments
\newtheorem{theorem}{Theorem}
\newtheorem{definition}[theorem]{Definition}
\newtheorem{corollary}[theorem]{Corollary}

% Shortcuts
\newcommand{\half}{\tfrac{1}{2}}
\newcommand{\third}{\tfrac{1}{3}}
\newcommand{\sixth}{\tfrac{1}{6}}
\newcommand{\twelfth}{\tfrac{1}{12}}
\newcommand{\CC}{\mathbb{C}}
\newcommand{\RR}{\mathbb{R}}
\newcommand{\HH}{\mathcal{H}}
\newcommand{\Vone}{V_1}
\newcommand{\Vhalf}{V_{1/2}}
\newcommand{\Vthird}{V_{1/3}}
\newcommand{\Pp}{P_+}
\newcommand{\Pm}{P_-}

\begin{document}

\preprint{UD/2026-001}

\title{Unified Dynamics: Spectral Geometry from Recursive Tetrahedral Transport}

\author{Benjamin Walker Mayes}
 \email{benjamin.w.mayes@gmail.com}
\affiliation{Unified Dynamics Alliance}

\date{\today}

\begin{abstract}
We construct a theory of physics from one geometric primitive: the tetrahedron. A doubly stochastic transport operator~$T$ on the six-dimensional edge space~$\CC^6$ of the tetrahedron yields eigenvalues $\{1,\half,\third\}$ with eigenspace dimensions $\{1,2,3\}$, matching the dimensional structure of $\mathrm{U}(1)\times\mathrm{SU}(2)\times\mathrm{SU}(3)$. We prove that the tetrahedron is the \emph{unique} simplex producing this decomposition. The Hodge star on the edge space splits $\CC^6$ into chiral sectors with dimensions $(2,1,1,2)$, yielding exact chirality without postulation. A Yukawa coupling matrix constructed from the tetrahedral geometry has singular values $1/\sqrt{3}$ and $1/3$ with basis-independent ratio~$\sqrt{3}$. The theory produces zero-parameter predictions for the fine-structure constant ($\alpha=7/960$, within $0.03\%$ of experiment), the CMB spectral tilt ($n_s=185/192$, within $0.14\%$), and lepton masses (within $0.8\%$). Two predictions currently fail: the Higgs mass ratio ($2.5\times$ discrepancy) and $V_{ub}$ ($2.5\times$ overshoot). The entire construction uses zero free parameters and one geometric input.
\end{abstract}

\maketitle

\section{\label{sec:intro}Introduction}

The Standard Model of particle physics contains 19~free parameters whose values are determined by experiment, not by the theory itself.\cite{PDG2024} Grand unified theories reduce this number but introduce new parameters at higher scales. A natural question persists: \emph{is there a geometric structure from which the gauge group, particle content, and coupling hierarchy follow with zero free parameters?}

We construct such a structure from one primitive: the tetrahedron. The construction proceeds in five steps:
\begin{enumerate}
\item A doubly stochastic transport operator~$T$ on the tetrahedral edge space~$\CC^6$ yields a rigid spectral decomposition (Sec.~\ref{sec:substrate}).
\item The spectrum $\{1,\half,\third\}$ with dimensions $\{1,2,3\}$ is unique among all simplices (Sec.~\ref{sec:uniqueness}).
\item Chirality, particle sectors, and massless gauge modes emerge from the Hodge star and eigenspace degeneracy (Sec.~\ref{sec:particles}).
\item A Yukawa coupling hierarchy with ratio~$\sqrt{3}$ and a $60^\circ$ coupling floor follow algebraically (Sec.~\ref{sec:mass}).
\item CKM-like mixing arises from inter-generation shell deformation (Sec.~\ref{sec:mixing}).
\end{enumerate}
The theory produces 15~falsifiable predictions (Sec.~\ref{sec:predictions}), of which 5~are exact structural matches, 8~are numerical matches within $20\%$, and 2~are honest failures.


\section{\label{sec:substrate}The Tetrahedral Substrate}

\subsection{\label{ssec:transport}Transport operator}

Consider the tetrahedron with vertices $\{0,1,2,3\}$ and edges~$\sigma_d$ labeled by the six 2-element subsets of $\{0,1,2,3\}$. Define the edge space $\HH=\CC^6$ with basis $\{|\sigma_d\rangle\}_{d=1}^{6}$.

The transport operator $T:\HH\to\HH$ is the $6\times6$ doubly stochastic matrix defined by the Sierpi\'nski iterated function system (IFS) weights on the tetrahedral attractor:
\begin{equation}\label{eq:transport}
T_{ij}=\begin{cases}
\half & i=j,\\[4pt]
\sixth & \sigma_i,\sigma_j\text{ opposite},\\[4pt]
\twelfth & \text{otherwise}.
\end{cases}
\end{equation}
Column and row sums each equal unity.

\subsection{\label{ssec:spectrum}Spectral decomposition}

The eigenvalues of~$T$ are
\begin{equation}\label{eq:eigenvalues}
\lambda_1=1,\quad \lambda_2=\half,\quad \lambda_3=\third,
\end{equation}
with eigenspace dimensions
\begin{equation}\label{eq:dimensions}
\dim\Vone=1,\quad \dim\Vhalf=2,\quad \dim\Vthird=3.
\end{equation}
These follow from the irreducible decomposition of~$\CC^6$ under the symmetric group~$S_4$:
\begin{equation}\label{eq:decomposition}
\CC^6 = \Vone\oplus\Vhalf\oplus\Vthird,
\end{equation}
where $\Vone$~carries the trivial representation, $\Vhalf$~the two-dimensional irrep, and $\Vthird$~the standard (three-dimensional) representation of~$S_4$.

\subsection{\label{ssec:commutant}Commutant algebra}

The algebra of operators commuting with~$T$ decomposes as
\begin{equation}\label{eq:commutant}
\mathrm{End}_{S_4}(\CC^6)\cong\mathfrak{gl}(1)\oplus\mathfrak{gl}(2)\oplus\mathfrak{gl}(3).
\end{equation}
This is the Standard Model gauge algebra skeleton: $\mathrm{U}(1)\times\mathrm{SU}(2)\times\mathrm{SU}(3)$.

\subsection{\label{ssec:hodge}Spin and Hodge structure}

The \emph{spin operator} $S:\HH\to\HH$ is defined by $S|\sigma_d\rangle=|\bar{\sigma}_d\rangle$, where $\bar{\sigma}_d$ is the complementary edge (the unique edge sharing no vertices with~$\sigma_d$). Then $S^2=I$, $[T,S]=0$, and $S$ decomposes~$\HH$ into $\pm1$ eigenspaces $\Pp$, $\Pm$ with $\dim\Pp=\dim\Pm=3$.

The Hodge star~$*$ on $\bigwedge^2\RR^4$ restricted to the edge space satisfies $*^2=I$ and $[S,*]=0$. The joint $(S,*)$ decomposition yields four chiral sectors:
\begin{equation}\label{eq:chirality}
\HH=\underbrace{(2)_\mathrm{LH}}_{\Pp}\oplus\underbrace{(1)_\mathrm{RH}}_{\Pp}\oplus\underbrace{(1)_\mathrm{LH}}_{\Pm}\oplus\underbrace{(2)_\mathrm{RH}}_{\Pm}.
\end{equation}
The chiral structure $(2,1,1,2)$ is exact and geometric.


\section{\label{sec:uniqueness}Tetrahedral Uniqueness}

\begin{theorem}[Tetrahedral Uniqueness]\label{thm:uniqueness}
Among all $n$-simplices ($n\geq2$), the tetrahedron ($n=3$) is the unique case whose edge representation under~$S_{n+1}$ decomposes into irreps of dimensions $\{1,2,3\}$.
\end{theorem}

\begin{proof}
The permutation representation of~$S_{n+1}$ on 2-element subsets of $[n{+}1]$ decomposes as\cite{FultonHarris}
\begin{equation}\label{eq:simplex_decomp}
\CC^{\binom{n+1}{2}}=V_{(n+1)}\oplus V_{(n,1)}\oplus V_{(n-1,2)},
\end{equation}
with dimensions $1$, $n$, and $n(n{-}1)/2-1$ respectively. Setting $\{1,\,n,\,n(n{-}1)/2-1\}=\{1,2,3\}$ yields $n=3$ as the unique solution:
\begin{itemize}
\item $n=2$: dimensions $\{1,2,0\}$---fails (zero-dimensional component).
\item $n=3$: dimensions $\{1,3,2\}=\{1,2,3\}$---succeeds.
\item $n\geq4$: $n\geq4$ and $n(n{-}1)/2-1\geq5$---neither equals 2 or 3.
\end{itemize}
Verified computationally for $n=2$ through $n=8$. Non-simplicial geometries (octahedron: 4~irreps; cube: 4; icosahedron: 7) all fail.
\end{proof}


\section{\label{sec:spine}The Loop Spine}

The spectral decomposition~(\ref{eq:decomposition}) generates a hierarchy of dynamical behaviors.

\begin{definition}[Transport, Closure, Reconstruction]\label{def:TCR}
For $|\psi\rangle\in\HH$, define three nested loop classes:
\begin{align}
L_T &: |\psi\rangle\neq0 \implies T^n|\psi\rangle\neq0\;\forall\,n, \label{eq:LT}\\
L_C &: T|\psi\rangle=\lambda|\psi\rangle\text{ for some }\lambda, \label{eq:LC}\\
L_R &: |\psi\rangle\in V_\lambda\text{ with }\dim V_\lambda\geq2. \label{eq:LR}
\end{align}
\end{definition}

These satisfy $L_R\subset L_C\subset L_T$.

\begin{theorem}[Transport Convergence]\label{thm:convergence}
For any $|\psi\rangle\in\HH$,
\begin{equation}\label{eq:convergence}
\bigl\|T^n-P_1\bigr\|=\Bigl(\frac{1}{2}\Bigr)^{\!n},
\end{equation}
where $P_1$ is the projector onto~$\Vone$.
\end{theorem}

\begin{proof}
$T$ is diagonalizable with eigenvalues $\{1,\half,\third\}$, so $T^n=P_1+(\half)^n P_{1/2}+(\third)^n P_{1/3}$. The operator norm of $T^n-P_1$ equals $(\half)^n$.
\end{proof}

\begin{theorem}[Repair Cost]\label{thm:repair}
The cost of transitioning a perturbation from~$V_\lambda$ to~$V_\mu$ is
\begin{equation}\label{eq:repair_cost}
C(\lambda\to\mu)=1-\mu.
\end{equation}
Cost depends only on the target eigenspace.
\end{theorem}

The \emph{repair burden} of a perturbation~$|\delta\psi\rangle$ is
\begin{equation}\label{eq:burden}
B(\delta\psi)=\sum_\mu(1-\mu)\,\|P_\mu|\delta\psi\rangle\|^2.
\end{equation}
In degenerate eigenspaces ($\dim V_\lambda\geq2$), within-eigenspace perturbations have $B=0$: reconstruction is free. This structurally protects massless gauge modes.


\section{\label{sec:particles}Particle Structure and Chirality}

\subsection{\label{ssec:ql}Quark-lepton decomposition}

The packet state space decomposes as
\begin{equation}\label{eq:ql_split}
\CC^8=\CC^6\oplus\CC^2,
\end{equation}
where $\CC^6$ (edges) supports the quark sector and $\CC^2$ (vertex excess) supports the lepton sector.

\subsection{\label{ssec:gluons}Massless gluons}

The color sector satisfies $T|_{\Vthird}=\third\cdot I_3$. Since $\Vthird$ is three-fold degenerate,
\begin{equation}\label{eq:massless}
B(\delta\psi)=0\quad\forall\;|\delta\psi\rangle\in\Vthird.
\end{equation}
Degeneracy protects the color sector from mass generation without a separate symmetry-breaking mechanism.

\subsection{\label{ssec:gen}Generation structure}

Three generations arise from dihedral stability of the IFS at levels $n=4,5,6$. Coherence retention per level:
\begin{equation}\label{eq:eta}
\eta=1-\frac{\Delta}{N}=\frac{185}{192}\approx0.964.
\end{equation}

\subsection{\label{ssec:corridor}Corridor selection dynamic}

The IFS corridor rotation~$U(\theta)$ respects the spin operator~$S$. Since $\Vhalf\subset\Pp$ and $\Vthird=\Pm$:
\begin{align}
U(\theta)\big|_{\Vone}&=I &&\text{(corridor-invariant: free)}, \label{eq:corr_V1}\\
U(\theta)\big|_{\Vthird}&=I &&\text{(corridor-invariant: confined)}, \label{eq:corr_V13}\\
U(\theta)\big|_{\Vhalf}&\neq I &&\text{(corridor-rotated: broken)}. \label{eq:corr_V12}
\end{align}
The physical role assignment---confinement for color, symmetry breaking for weak, freedom for gravity---is forced by the $(S,*)$ algebra.


\section{\label{sec:mass}Mass and Coupling Hierarchy}

\begin{theorem}[Yukawa Mass Ratio]\label{thm:yukawa}
The $2\times2$ Yukawa coupling matrix~$Y$ on the tetrahedral substrate has singular values $\sigma_1=1/\sqrt{3}$, $\sigma_2=1/3$. Their ratio is
\begin{equation}\label{eq:yukawa_ratio}
\frac{\sigma_1}{\sigma_2}=\sqrt{3}.
\end{equation}
Both $\det(Y)=1/(3\sqrt{3})$ and $\mathrm{tr}(Y^\dagger Y)=4/9$ are basis-independent. The ratio is invariant under all~$\Vhalf$ rotations.
\end{theorem}

\begin{theorem}[Coupling Angle Floor]\label{thm:angle}
The intra-generation coupling angle~$\theta$ satisfies
\begin{equation}\label{eq:angle_floor}
|\cos\theta|_\mathrm{max}=\frac{\sigma_1^2-\sigma_2^2}{\sigma_1^2+\sigma_2^2}=\frac{1}{2},
\end{equation}
giving $\theta\in[60^\circ,120^\circ]$. The $60^\circ$ floor is the tetrahedral angle $\arccos(\half)$.
\end{theorem}


\section{\label{sec:mixing}Inter-Generation Mixing}

The $60^\circ$ floor (Theorem~\ref{thm:angle}) is an \emph{intra}-generation result. CKM mixing is \emph{inter}-generation, controlled by the shell deformation~$\varepsilon(n)$:
\begin{equation}\label{eq:epsilon}
\varepsilon(4)=30^\circ,\quad\varepsilon(5)=12^\circ,\quad\varepsilon(6)=0^\circ.
\end{equation}
The Cabibbo angle:
\begin{equation}\label{eq:cabibbo}
\sin\theta_C=\sin(\Delta\varepsilon_{23})=\sin12^\circ=0.208,
\end{equation}
compared to experiment: $0.225$ (discrepancy $7.6\%$). The element~$|V_{cb}|$:
\begin{equation}\label{eq:vcb}
|V_{cb}|\approx\lambda^2=0.043,
\end{equation}
compared to experiment: $0.041$ (discrepancy $5.4\%$).\cite{Wolfenstein} A CP phase $\delta=\pi/30=6^\circ$ arises structurally from $\varepsilon(4)-2\varepsilon(5)$.

The element $|V_{ub}|\approx\lambda^3=0.009$ overshoots experiment ($0.0036$) by $2.5\times$. Resolution requires the full three-generation Yukawa with the CP phase, which remains open.


\section{\label{sec:predictions}Predictions}

All predictions below follow from one geometric input with zero free parameters.

\subsection{Exact structural matches}

\begin{table}[h]
\caption{\label{tab:exact}Exact structural predictions.}
\begin{ruledtabular}
\begin{tabular}{lcc}
Quantity & UD & Experiment\\
\hline
Spatial dimensions $d$ & 3 & 3\\
Generations $N_\mathrm{gen}$ & 3 & 3\\
Colors $N_c$ & 3 & 3\\
Gluons & 8 & 8\\
Graviton polarizations & 2 & 2\\
\end{tabular}
\end{ruledtabular}
\end{table}

\subsection{Numerical predictions}

\begin{table}[h]
\caption{\label{tab:numerical}Numerical predictions.}
\begin{ruledtabular}
\begin{tabular}{lccc}
Quantity & UD & Experiment & Error\\
\hline
$\alpha$ & $7/960$ & $1/137.036$ & $0.03\%$\\
$n_s$ & $185/192$ & $0.9649\pm0.0042$ & $0.14\%$\\
$m_\mu$ (MeV) & $105.9$ & $105.66$ & $0.2\%$\\
$m_\tau$ (MeV) & $1762$ & $1776.9$ & $0.8\%$\\
Cabibbo angle & $\sin12^\circ$ & $0.225$ & $7.6\%$\\
$|V_{cb}|$ & $0.043$ & $0.041$ & $5.4\%$\\
$\sin^2\theta_W$ (bare) & $1/4$ & $0.231$ & $8.1\%$\\
$\sigma_1/\sigma_2$ & $\sqrt{3}$ & ${\sim}2.13$ & ${\sim}19\%$\\
\end{tabular}
\end{ruledtabular}
\end{table}

\subsection{Honest failures}

\begin{table}[h]
\caption{\label{tab:failures}Known failures.}
\begin{ruledtabular}
\begin{tabular}{lccc}
Quantity & UD & Experiment & Status\\
\hline
$m_H/m_W$ & 3.9 & 1.56 & Fails ($2.5\times$)\\
$|V_{ub}|$ & 0.009 & 0.0036 & Fails ($2.5\times$)\\
\end{tabular}
\end{ruledtabular}
\end{table}

Both failures have identified paths to resolution within the existing framework.

\subsection{Falsifiable tests}

\begin{enumerate}
\item Discovery of a 4th generation falsifies the dihedral stability argument.
\item Proton decay below $10^{45}$~yr falsifies the confinement mechanism.
\item Tree-level flavor-changing neutral currents falsify the corridor dynamic.
\item CMB-S4 measurement of~$n_s$ with $\sigma<0.002$ will confirm or exclude $185/192$.
\item RG running of $\sin^2\theta_W$ from $1/4$ to $M_Z$ will close or widen the $8\%$ gap.
\end{enumerate}


\section{\label{sec:discussion}Discussion}

The central result is that one geometric object---the tetrahedron---produces, through its transport operator and spectral decomposition, the dimensional structure of the Standard Model gauge group, exact chirality, a mass coupling hierarchy, and partial CKM recovery, all with zero free parameters.

Theorem~\ref{thm:uniqueness} answers ``why the tetrahedron?''\ with proof. The corridor selection dynamic (Sec.~\ref{ssec:corridor}) assigns physical roles: confinement for~$\Vthird$ ($\Pm$, corridor-invariant), symmetry breaking for~$\Vhalf$ ($\Pp$, corridor-rotated), freedom for~$\Vone$ (trivially invariant). This is forced by $(S,*)$, not postulated.

The theory's strength is constraint: no tuning is possible. Its weakness is two numerical failures (Higgs ratio, $V_{ub}$), representing bounded open problems with identified paths.


\section{\label{sec:conclusion}Conclusion}

Unified Dynamics proposes that gauge structure, particle content, and coupling hierarchy emerge from the spectral geometry of one recursive tetrahedral substrate with zero free parameters. The construction produces 13~predictions in agreement with experiment and 2~honest failures. Specific measurements---fourth generation, proton decay, tree-level FCNC, precision $n_s$---would falsify its central mechanisms.

If this construction is wrong, it is wrong in a specific language. If it is right, much of physics may be the spectrum of one small geometric fact.


\begin{acknowledgments}
The author thanks the Unified Dynamics Alliance for collaborative development and internal review of the theorem registry (583~locked theorems, April 2026).
\end{acknowledgments}


\section*{Data Availability Statement}

Data sharing is not applicable to this article as no new data were created or analyzed in this study. All results follow from the mathematical framework described herein.


\appendix

\section{\label{app:constants}Derived Constants}

The following constants are derived from the tetrahedral substrate with $\Delta=7$, $N=192$:
\begin{align}
x &= \frac{\Delta}{N}=\frac{7}{192}, \label{eq:x}\\
\alpha_0 &= \frac{7}{960}, \label{eq:alpha0}\\
n_s &= \frac{185}{192}=1-\frac{7}{192}, \label{eq:ns}\\
S_0 &= \frac{43}{7}, \label{eq:S0}\\
\eta &= 1-x = \frac{185}{192}. \label{eq:eta_app}
\end{align}


\nocite{*}
\bibliography{aipsamp}% Replace with your .bib file

\end{document}
% ****** End of file UD_Manuscript_APS.tex ******
